% Part: sets-functions-relations
% Chapter: sets
% Section: proofs-about-sets

\documentclass[../../../include/open-logic-section]{subfiles}

\begin{document}

\olfileid{sfr}{set}{prf}
\olsection{د سټونو په اړه ثبوتونه}

\begin{editorial}
د ثبوتونو د طريقو عمومي باب په \verb|content/methods/proofs| کښې دے۔ دلته پر سټونو د دغو طريقو يوه ځانګړې بېلګه لولو۔
\end{editorial}

\begin{explain}
سټونه او تر دې ځايه معرفي شوے نښليک د سټونو د ټاکلو او د هغو تر منځ د اړيکو د څرګندولو لنډه لاره ده۔ ډېر وخت به د دغو اړيکو په باب دعوې هم ثابتوو۔ که له رياضيکي ثبوتونو سره اشنا نۀ يئ، يوه ساده بېلګه ګورو۔ د هر دوو سټونو $X$ او $Y$ دپاره ښيو چې تل $X \cap
(X \cup Y) = X$۔ د سټونو داسې برابري څنګه ثابتوو؟ په ياد ولرئ چې سټونه يوازې له خپلو !!{element}s ټاکل کېږي؛ دوه سټونه هغه او يوازې هغه وخت برابر دي چې هماغه !!{element}s ولري۔ نو دلته ثابتوو چې (a) د $X \cap (X \cup
Y)$ هر !!{element} د $X$ هم !!{element} دے، او برعکس (b) د $X$ هر !!{element} د $X \cap (X \cup Y)$ هم !!{element} دے۔ يعنې دواړه ښيو، (a) $X \cap (X \cup Y) \subseteq X$ او (b) $X \subseteq X \cap (X \cup Y)$۔

د «د $X \cap (X
\cup Y)$ هر !!{element}~$z$ د $X$ هم !!{element} دے» په څېر د يوې عامې دعوې د ثابتولو دپاره لومړے فرضوو چې $z \in X \cap (X \cup Y)$ يو اختياري غړے راکړے شوے او له همدې فرضه $z \in X$ ثابتوو۔ دا د «عمومي شرطي ثبوت» په نوم هم پېژندلے شئ۔ د فرض او پايلې تعريفونه هم کاروو، لکه دلته چې د $\cap$ او $\cup$ تعريفونه کاروو۔ نو د (a) دليل داسې دے۔

\begin{quote}
(a) لومړے غواړو وښيو چې $X \cap (X \cup Y) \subseteq X$۔ د $\subseteq$ د تعريف له مخې، د هر $z$ دپاره، که $z \in X \cap (X \cup Y)$ نو $z \in X$۔ فرض کړئ $z \in X \cap (X \cup
Y)$۔ څرنګه چې $z$ د دوو سټونو د تقاطع !!{element} هغه او يوازې هغه وخت دے چې د دواړو سټونو !!{element} وي، نو $z \in X$ او $z \in X \cup Y$ اخلو۔ په ځانګړي ډول $z \in X$، هماغه چې ثابتول مو غوښتل۔
\end{quote}

د ثبوت لومړۍ نيمه پوره شوه۔ په وروستي ګام کښې مو وکارول چې که د او تړنه، يعنې $z \in X$ او $z \in X
\cup Y$، له يوه فرضه راځي، نو هره تړلې برخه ئې له هماغه فرضه راځي۔ دغه قاعده «د او ايستل» يا $\land$ ايستل بولي۔ اوس (b) ثابتوو۔

\begin{quote}
(b) اوس ثابتوو چې $X \subseteq X \cap (X \cup Y)$۔ د $\subseteq$ د تعريف له مخې، د هر $z$ دپاره که $z \in X$، نو $z \in X \cap
(X \cup Y)$ هم۔ فرض کړئ $z \in X$۔ د $z \in X
\cap (X \cup Y)$ د ښودلو دپاره د $\cap$ له تعريفه دوه خبرې غواړو، (i) $z \in X$ او (ii) $z \in X \cup Y$۔ (i) پخپله زمونږ فرض دے، نو نور ثبوت نۀ غواړي۔ د (ii) دپاره په ياد ولرئ چې $z$ د سټونو د اتحاد !!{element} هغه او يوازې هغه وخت دے چې لږ تر لږه د يوه سټ !!{element} وي۔ دلته دا شرط پوره دے، ځکه $z \in X$ او $X \cup Y$ د $X$ او $Y$ اتحاد دے۔ نو $z \in X \cup Y$۔ (i) $z \in X$ او (ii) $z \in X \cup Y$ دواړه مو وښودل، نو د $\cap$ د تعريف له مخې $z \in X \cap (X \cup Y)$۔
\end{quote}

دا بيان لږ اوږد دے، خو ښيي چې د سټونو او د هغو د اړيکو په اړه څنګه استدلال کوو۔ عموماً دومره تفصيل نۀ ورکوو؛ په ځانګړي ډول ښايي ټول تعريفونه بيا وانۀ يوو۔ په يوه پرمخ تللي متن کښې همدا ثبوت ډېر لنډ داسې ليکلے شو۔
\end{explain}

\begin{prop}[جذب]
د هر دوو سټونو $X$، $Y$ دپاره،
\[
X \cap (X \cup Y) = X
\]
\end{prop}

\begin{proof}
(a) فرض کړئ $z \in X \cap (X \cup Y)$۔ نو $z \in X$، له دې امله $X \cap (X
  \cup Y) \subseteq X$۔

(b) اوس فرض کړئ $z \in X$۔ نو $z \in X \cup Y$ هم، او له دې امله $z \in X \cap (X \cup Y)$ هم۔ په نتيجه کښې $X \subseteq X \cap (X \cup
  Y)$۔
\end{proof}

\begin{prob}
په تفصيل ثابت کړئ چې $X \cup (X \cap Y) = X$۔ بيا يو لنډ ثبوت ورکړئ۔ (اشاره۔ د $X \cup (X \cap Y) \subseteq X$ لوري دپاره په جلا حالتونو ثبوت، يعنې $\lor$ ايستل، پکار دي۔)
\end{prob}

\end{document}
